% !TeX program = xelatex
\documentclass[a4paper,10pt]{article}

% ============================================================================
% ក្រដាសដាច់ដោយឡែក៖ សេណារីយ៉ូនៃការបង្រៀន (Scenario នៃការបង្រៀន – ៥០ នាទី)
% ============================================================================
\usepackage[a4paper, left=0.95cm, right=0.95cm, top=0.88cm, bottom=0.98cm, footskip=0.50cm]{geometry}
\usepackage{amsmath}
\usepackage{unicode-math}
\usepackage{fontspec}
\usepackage{array, longtable, colortbl}
\usepackage[table]{xcolor}
\usepackage{pifont}
\usepackage{fancyhdr}
\usepackage{tikz}

% LaTeX 2025 + longtable compatibility guard
\makeatletter
\let\@expl@@@mark@update@singlecol@structures@@\relax
\makeatother

% Cross-platform font path resolver (Windows local + Linux/Cloud repo fonts/)
\providecommand{\NGSFindFonts}{%
  \IfFileExists{C:/Windows/Fonts/cambria.ttc}{%
    \def\NGSSysPath{C:/Windows/Fonts/}%
    \def\NGSUserPath{C:/Users/KIT/AppData/Local/Microsoft/Windows/Fonts/}%
  }{%
    \IfFileExists{./fonts/KhmerOSsiemreap.ttf}{\def\NGSSysPath{./fonts/}\def\NGSUserPath{./fonts/}}{%
    \IfFileExists{../fonts/KhmerOSsiemreap.ttf}{\def\NGSSysPath{../fonts/}\def\NGSUserPath{../fonts/}}{%
    \IfFileExists{../../fonts/KhmerOSsiemreap.ttf}{\def\NGSSysPath{../../fonts/}\def\NGSUserPath{../../fonts/}}{%
    \IfFileExists{../../../fonts/KhmerOSsiemreap.ttf}{\def\NGSSysPath{../../../fonts/}\def\NGSUserPath{../../../fonts/}}{%
    \IfFileExists{../../../../fonts/KhmerOSsiemreap.ttf}{\def\NGSSysPath{../../../../fonts/}\def\NGSUserPath{../../../../fonts/}}{%
    \IfFileExists{../../../../../fonts/KhmerOSsiemreap.ttf}{\def\NGSSysPath{../../../../../fonts/}\def\NGSUserPath{../../../../../fonts/}}{%
    \IfFileExists{../../../../../../fonts/KhmerOSsiemreap.ttf}{\def\NGSSysPath{../../../../../../fonts/}\def\NGSUserPath{../../../../../../fonts/}}{%
      \def\NGSSysPath{../../../../../../../fonts/}\def\NGSUserPath{../../../../../../../fonts/}%
    }}}}}}}%
  }%
}
\NGSFindFonts

\setmathfont[
  Path={\NGSSysPath},
  FontIndex=1,
  Scale=1.02
]{cambria.ttc}

\defaultfontfeatures{Ligatures=TeX}

\setmainfont[
  Path={\NGSUserPath},
  Extension=.ttf,
  Script=Khmer,
  BoldFont={KhmerOSsiemreap},
  BoldFeatures={FakeBold=1.8},
  ItalicFont={KhmerOSsiemreap},
  ItalicFeatures={FakeSlant=0.25},
  BoldItalicFont={KhmerOSsiemreap},
  BoldItalicFeatures={FakeBold=1.8, FakeSlant=0.25}
]{KhmerOSsiemreap}

\newfontfamily\muol[
  Path={\NGSUserPath},
  Extension=.ttf,
  Script=Khmer
]{KhmerOSmuollight}

\newfontfamily\TimesEn[
  Path={\NGSSysPath},
  Extension=.ttf,
  Scale=1.00,
  BoldFont={timesbd},
  ItalicFont={timesi},
  BoldItalicFont={timesbi}
]{times}

\newfontfamily\NagaBold[
  Path={\NGSUserPath},
  Extension=.ttf
]{AKbalthom Naga Bold}

\newfontfamily\CalibriFont[
  Path={\NGSSysPath},
  Extension=.ttf
]{calibri}

\XeTeXlinebreaklocale "km"
\XeTeXlinebreakskip = 0pt plus 1pt minus 0.1pt

\definecolor{stepgray}{HTML}{D9D9D9}
\definecolor{headgray}{HTML}{F2F2F2}
\definecolor{darkred}{HTML}{C00000}
\definecolor{red5e}{HTML}{FF0000}

\pagestyle{fancy}
\fancyhf{}
\renewcommand{\headrulewidth}{0pt}
\fancyfoot[R]{\CalibriFont\fontsize{11pt}{13pt}\selectfont\thepage}

\setlength{\parindent}{0pt}
\setlength{\parskip}{0pt}
\setlength{\tabcolsep}{4pt}

\newcolumntype{L}[1]{>{\raggedright\arraybackslash}p{#1}}
\newcommand{\en}[1]{{\TimesEn #1}}
\newcommand{\enbf}[1]{{\TimesEn\textbf{#1}}}
\newcommand{\fiveE}[1]{{\NagaBold\fontsize{9pt}{12.5pt}\selectfont\textcolor{red5e}{#1}}}
\newcommand{\ditem}{\tikz[baseline=-0.5ex]\fill (0,0.40ex) -- (0.40ex,0) -- (0,-0.40ex) -- (-0.40ex,0) -- cycle;\hspace{0.25em}}
\newcommand{\chkbox}{\tikz[baseline=-0.1ex]\draw[black, line width=0.75pt](0,0) rectangle (0.23,0.23);}
\newcommand{\cnum}[1]{{\fontsize{9.5pt}{11pt}\selectfont\ding{\numexpr181+#1\relax}}\hspace{0.15em}}
\newcommand{\tri}[1]{\ensuremath{\triangle #1}}
\newcommand{\ang}[1]{\ensuremath{\angle #1}}
\newcommand{\cg}{\ensuremath{\cong}}
\newcommand{\ncg}{\ensuremath{\not\cong}}

\begin{document}
\fontsize{10pt}{14.8pt}\selectfont

% -------------------- ក្បាលឯកសារ Scenario --------------------
\begin{center}
  {\muol\fontsize{11pt}{16pt}\selectfont ព្រះរាជាណាចក្រកម្ពុជា}\\[1pt]
  {\muol\fontsize{10pt}{15pt}\selectfont ជាតិ សាសនា ព្រះមហាក្សត្រ}\\[4pt]
  {\muol\fontsize{10.5pt}{15pt}\selectfont សេណារីយ៉ូនៃការបង្រៀន \enbf{(Scenario} នៃការបង្រៀន\enbf{)}}
\end{center}

\vspace{2pt}
{\renewcommand{\arraystretch}{1.06}
\begin{tabular}{@{}l l l p{0.5cm} l l l@{}}
  - អនុវិទ្យាល័យ & : & ព្រះក្រសាំង (កំពង់ស្វាយ-កំពង់ធំ) & & - មុខវិជ្ជា & : & គណិតវិទ្យា ថ្នាក់ទី៨ \\
  - មេរៀនទី១២ & : & ១.៣.២ អ.ជ (HL) និងប្រសព្វពុះមុំក្នុង & & - រយៈពេលសរុប & : & \textcolor{darkred}{\textbf{១ ម៉ោងសិក្សា (៥០ នាទី)}} \\
  - គ្រូបង្រៀន & : & ហុក ប៊ុនថែន & & - កាលបរិច្ឆេទ & : & ថ្ងៃចន្ទ ទី៣០ ខែវិច្ឆិកា ឆ្នាំ២០២៦ \\
\end{tabular}
}

\vspace{3pt}

% -------------------- I. តារាងបែងចែកពេលវេលា និងសេណារីយ៉ូលម្អិត (៥០ នាទី) --------------------
\textbf{I.} {\muol តារាងបែងចែកពេលវេលា និងសេណារីយ៉ូប្រតិបត្តិការក្នុងថ្នាក់} \textbf{(៥០ នាទី)}

\renewcommand{\arraystretch}{1.10}
\begin{longtable}{|c|L{2.75cm}|c|L{7.65cm}|L{4.45cm}|}
\hline
\rowcolor{headgray}
\multicolumn{1}{|c|}{\muol នាទី} &
\multicolumn{1}{c|}{{\muol ជំហាន \enbf{/ 5Es}}} &
\multicolumn{1}{c|}{\muol រយៈពេល} &
\multicolumn{1}{c|}{{\muol សេណារីយ៉ូប្រតិបត្តិ \enbf{(}ពាក្យពេចន៍គ្រូ \enbf{\&} សម្ភារៈ\enbf{)}}} &
\multicolumn{1}{c|}{{\muol ការគ្រប់គ្រង \enbf{\&} វិធីបម្រុង}} \\
\hline
\endfirsthead

\hline
\rowcolor{headgray}
\multicolumn{1}{|c|}{\muol នាទី} &
\multicolumn{1}{c|}{{\muol ជំហាន \enbf{/ 5Es}}} &
\multicolumn{1}{c|}{\muol រយៈពេល} &
\multicolumn{1}{c|}{{\muol សេណារីយ៉ូប្រតិបត្តិ \enbf{(}ពាក្យពេចន៍គ្រូ \enbf{\&} សម្ភារៈ\enbf{)}}} &
\multicolumn{1}{c|}{{\muol ការគ្រប់គ្រង \enbf{\&} វិធីបម្រុង}} \\
\hline
\endhead

\hline
\endfoot

\hline
\endlastfoot

00--02 & \textbf{ជំហានទី១}\newline រដ្ឋបាលថ្នាក់ & ២ នាទី &
\ditem ពិនិត្យអវត្តមាន និងអនាម័យថ្នាក់រៀន ។ &
ចាត់ប្រធានក្រុមស្រាប់ ។ \\[3pt]
& & &
\ditem ចាត់ក្រុម ៣ នាក់ («ក្រុមវិស្វករសួនរង្វង់ត្រីកោណ») ចែកសន្លឹក ១:១ (ទំព័រ ១) ដាក់ផ្កាប់ទុកលើតុ ។ &
ហាមកាត់ក្រដាសមុនបញ្ជាគ្រូ ។ \\ \hline

02--05 & \textbf{ជំហានទី២ (ក)}\newline កំណែកិច្ចការផ្ទះ & ៣ នាទី &
\ditem \textbf{សំណួរបញ្ឆោត (\#6 ទំ.១៦៤)}៖ «ក្នុងលំហាត់ទី៦ តើ $AD$ ជាអ៊ីប៉ូតេនុស ឬជាជ្រុងមុំកែង?» &
ផ្ដោតលើជ្រុងកែង vs អ៊ីប៉ូតេនុស ។ \\[3pt]
& & &
\ditem ផ្ទៀងខ្លី \#6 ($AD$ ជ្រុងកែងរួម $\to$ ម.ជ.ម) និង \#7 (ក. $PB=PC$ ជ.ម.ជ; ខ. $PH=PK$ អ.ម) ។ &
កុំឱ្យសិស្សសរសេរវែងលើក្ដារខៀន ។ \\ \hline

05--08 & \textbf{ជំហានទី២ (ខ)}\newline រំលឹកមេរៀនទាក់ទង & ៣ នាទី &
\ditem \textbf{ភាគ ១ (១.៥ ន.)}៖ រំលឹក $\ang{BAC}+\ang{CAB'}=90^\circ{+}90^\circ=180^\circ \Rightarrow B, A, B'$ ត្រង់ជួរ ($BB'$) និង $\tri{BCB'}$ សមបាទ ($CB=CB'\Rightarrow\ang{B}=\ang{B'}$) ។ &
គូសរូបភ្ជាប់ជ្រុង $AC$ លើក្ដារខៀនឱ្យសិស្សឃើញបន្ទាត់ត្រង់ $BB'$ ច្បាស់ ។ \\[3pt]
& & &
\ditem \textbf{ភាគ ២ (១.៥ ន.)}៖ រំលឹកលក្ខខណ្ឌ អ.ម (HA) និងលក្ខណៈផ្ទាល់ពុះមុំ ($MA=MB$) ។ &
រំលឹកអ៊ីប៉ូតេនុសឈម $90^\circ$ ។ \\ \hline

08--11 & \fiveE{E1 Engage}\newline សួនរង្វង់ ៣ ផ្លូវ & ៣ នាទី &
\ditem លើករឿង \textbf{«អាងទឹករង្វង់សួនផ្លូវប្រសព្វ ៣ ខ្សែ និងអាថ៌កំបាំងរួចផុតពីអន្ទាក់ ជ.ជ.ម»}៖ ពុះមុំ $\ang{B}, \ang{C}$ ប្រសព្វត្រង់ $I$, ផ្លូវកែង $ID, IE, IF$ ។ &
បង្ហាញរូប E1 លើក្ដារខៀន ។ \\[3pt]
& & &
\ditem ស្រង់មតិហេតុអ្វី $ID=IE=IF$ ហើយ $AI$ ពុះ $\ang{A}$ ទាំងមាន $90^\circ$ (មិននៅចន្លោះ) រួចសរសេរចំណងជើង ១.៣.២ ។ &
កុំទាន់ផ្ដល់ចម្លើយផ្លូវការ ទុកជាបេសកកម្មឱ្យសិស្សស្រាយក្នុង E4--E5 ។ \\ \hline

11--22 & \fiveE{E2 Explore}\newline ពិសោធន៍ ១ \& ២ & ១១ នាទី &
\ditem \textbf{សកម្មភាព ១ (៦ ន.)}៖ លើទំព័រ ១ (១:១) កាត់ត្រីកោណកែង $3$-$4$-$5\text{cm}$ ពីរ ($AC=A'C'=3\text{cm}, BC=B'C'=5\text{cm}$) ត្រឡប់ផ្គុំ $A'C'$ លើ $AC$ ($\ang{BAB'}=180^\circ$) ។ &
បោះពុម្ពទំព័រ ១ ខ្នាត \textbf{Actual Size 100\%} ។ \\[3pt]
& & &
+ ធៀបនឹងធ្នូ $35^\circ$ ($PQ=4\text{cm}, r=2.8\text{cm}$) កាត់កាំ $PX$ ២ កន្លែង ($R_1, R_2$) ទល់កែង ១ គត់ ($H$) ។ &
បង្ហាញអន្ទាក់ ជ.ជ.ម មុនបញ្ជាក់ អ.ជ ។ \\[3pt]
& & &
\ditem \textbf{សកម្មភាព ២ POE (៥ ន.)}៖ បត់ក្រដាស $\tri{ABC}$ សាមញ្ញ ($BC=9.2, AB=6.5, AC=8.0\text{cm}$) រកផ្នត់ពុះ $\ang{B}, \ang{C}$ ប្រសព្វ $I$, តេស្តផ្នត់ $\ang{A}$, រួចវាស់ $ID=IE=IF\approx 2.1\text{cm}$ (រង្វង់ចារឹកក្នុង) ធៀប $IA\approx 3.4, IB\approx 4.4, IC\approx 5.8\text{cm}$ ។ &
គ្រូដើរត្រួតផ្នត់ពុះមុំកុំឱ្យសិស្សបត់ជាមេដ្យាន ។ \\ \hline

22--28 & \fiveE{E3 Explain}\newline សម្រាយ \& \cnum{1}--\cnum{5} & ៦ នាទី &
\ditem ដឹកនាំសម្រាយត្រឡប់ $AC=A'C'$ ($\ang{BAB'}=180^\circ, \ang{B}=\ang{B'}$) និង\textbf{កែទំព័រ ១៥៤ ទាំង ៣ ចំណុច}៖ (១) $AC=A'C'$ ស៊ីគ្នា, (២) កែ \textit{«(តាមលកខណ្ឌ អ.ជ)»} $\to$ \textbf{«(តាមលក្ខខណ្ឌ អ.ម)»}, (៣) \textit{«លកខណ្ឌ»} $\to$ \textbf{«លក្ខខណ្ឌ»} ។ &
សង្កត់ធ្ងន់៖ ប្រើ អ.ជ ខ្លួនឯងមកស្រាយ អ.ជ គឺជារង្វង់តក្ក ត្រូវប្រើ អ.ម! \\[3pt]
& & &
\ditem បង្គុំទ្រឹស្ដីបទ \textbf{អ.ជ (HL)} (ទំ.១៥៥) + \textbf{ប្រសព្វពុះមុំក្នុង} (ទំ.១៥៦) + តារាង H--C--C និងបិទផ្ទាំង \textbf{ធនាគារហេតុផល \cnum{1}--\cnum{5}} (ទំព័រ ៤) ។ &
ញែក $I$ (ពុះមុំ) ពី $O$ និង $G$ ។ \\ \hline

28--38 & \fiveE{E4 Elaborate}\newline ច្រវាក់ អ.ម$\to$អ.ជ & ១០ នាទី &
\ditem ចែកកាត \textbf{«ផែនទីច្រវាក់ អ.ម $\to$ អ.ជ ផ្ចិត I»} (ទំព័រ ២, ទំ.១៥៥)៖\newline
+ ២.៥ន. \textbf{កែរូបទំ.១៥៥} (គូស $\tri{ABC}$ សាមញ្ញ, $AI$ ដាច់ៗ, \textbf{ដក \textbf{/} ពី $AD, BD, BE, AF$ និងធ្នូ $\ang{A}$}) រួច \textbf{[tracing]} $90^\circ\to BI, CI\to AI$ &
ព្រមានមុនសិស្សសរសេរ៖ ហាមយកសញ្ញា $/$ លើរូបទំ.១៥៥ ទៅប្រើ SAS ឬ SSS ដាច់ខាត! \\[3pt]
& & &
+ ៥ន. សរសេរសម្រាយ ២ ដំណាក់កាល៖ ក. អ.ម $\times 2 \Rightarrow ID=IE=IF$; ខ. អ.ជ ($\tri{IDA}\cg\tri{IFA}$) $\Rightarrow AI$ ពុះ $\ang{BAC}$ (\cnum{1}--\cnum{5}) &
បំពេញខ្សែសង្វាក់ អ.ម $\to$ អ.ជ តាមលំដាប់ ។ \\[3pt]
& & &
+ ២.៥ន. \textbf{EIA 30s} ញែក $ID=IE=IF$ (ចម្ងាយជ្រុង) ពី $IA, IB, IC$ (កំពូល) ។ &
ផ្ទៀងរង្វង់ចារឹកក្នុង vs ចារឹកក្រៅ ។ \\ \hline

38--48 & \fiveE{E5 Evaluate}\newline តុសវនកម្ម \& កាត & ១០ នាទី &
\ditem \textbf{«តុសវនកម្មចំណុចកណ្ដាល--អង្កត់កែង ៣ តួនាទី» (៧ ន.)} (ទំព័រ ៣, ប្រតិបត្តិទំ.១៥៦ $\to$ បំបែក $\tri{MDB}, \tri{MEC}$)៖\newline
+ \textbf{តួនាទី ១ (Pause-Ask)}៖ ពន្យល់មុនគូស $90^\circ\to$ អ៊ីប៉ូ $MB=MC$ ($M$ ចំណុចកណ្ដាល) &
កុំឱ្យសិស្សមើលភ្នែកថា $AB=AC$ ពីរូបទំ.១៥៦! \\[3pt]
& & &
+ \textbf{តួនាទី ២}៖ សរសេរ \cnum{1}--\cnum{4} ($\tri{MDB}\cg\tri{MEC}$ តាម អ.ជ)\newline
+ \textbf{តួនាទី ៣ (Pause-Review)}៖ ផ្ទៀង $\ang{B}=\ang{C}$ (\cnum{5}) + \textbf{បិទបញ្ហា E1} ។ &
បង្ហាញថា $AI$ ពុះ $\ang{A}$ ដោយ អ.ជ ។ \\[3pt]
& & &
\ditem \textbf{កាតត្រួតពិនិត្យ A--D (៣ ន. / 60--90s)} (ទំព័រ ៥)៖ (A) \ding{51} អ.ជ បង្វិល, (B) \ding{55} អន្ទាក់ $60^\circ$ ជ.ជ.ម ($4, 5\text{cm}$), (C) \ding{55} $IA=IB=IC$, (D) \ding{55} អ៊ីប៉ូ vs ជ្រុង $5\text{cm}$ + ម្រាម ១--៣ ។ &
បើអស់ពេល កាត់កាត A--D មក ៩០ វិនាទី ។ \\ \hline

48--50 & \textbf{ជំហានទី៥}\newline កិច្ចការផ្ទះ \& បណ្ដាំ & ២ នាទី &
\ditem ចែកប័ណ្ណកិច្ចការផ្ទះ (ទំព័រ ៦)៖ \#10 (ទំ.១៦៥ បំពេញ \textbf{[AC និង AB]} បំបែក $\tri{BDC}\cg\tri{CEB}$ តាម អ.ជ) \& \#12 (ទំ.១៦៥ $\tri{ACB}\cg\tri{BDA}$ តាម អ.ជ $\Rightarrow BC=AD, BC\parallel AD$) ។ &
បញ្ជូនរូបលំហាត់ចូល \en{Telegram / Classroom} ជាស្រេច ។ \\[3pt]
& & &
\ditem ផ្ដាំផ្ញើសុវត្ថិភាពចរាចរណ៍ និងជួយការងារផ្ទះ ។ &
បញ្ចប់ម៉ោងសិក្សាឱ្យទាន់ ៥០ នាទី ។ \\ \hline

\rowcolor{stepgray}
\multicolumn{2}{|r|}{\textbf{សរុបរយៈពេលបង្រៀនពេញ}} & \textbf{៥០ នាទី} &
\multicolumn{2}{l|}{ត្រូវគ្នា ១០០\% នឹង \en{LS\_Plan\_HL.tex} និង \en{Printable\_Materials\_HL.tex} (ទំព័រ ១--៦)} \\ \hline
\end{longtable}

\vspace{1pt}

% -------------------- II. បញ្ជីត្រៀមសម្ភារៈ & III. ការគ្រប់គ្រងកំហុស --------------------
{\fontsize{9.2pt}{13.0pt}\selectfont\renewcommand{\arraystretch}{1.04}
\begin{minipage}[t]{0.44\textwidth}
  \textbf{II.} {\muol បញ្ជីត្រៀមសម្ភារៈមុនចូលថ្នាក់} \textbf{(ស-ខ្មៅ)}\\[2pt]
  \begin{tabular}{|c|L{6.75cm}|}
    \hline
    \rowcolor{headgray}
    \textbf{\ding{51}} & \multicolumn{1}{c|}{\textbf{សម្ភារៈត្រូវរៀបចំទុកជាមុន}} \\ \hline
    \chkbox & សន្លឹក ១:១ ទំព័រ ១ (\textbf{Actual Size 100\%})៖ ត្រីកោណ $3$-$4$-$5\text{cm}$ ២ សន្លឹក, ធ្នូ $35^\circ$, ក្រដាសបត់ $\tri{ABC}$ សាមញ្ញ ($9.2\text{-}6.5\text{-}8.0\text{cm}$) \\ \hline
    \chkbox & កន្ត្រៃ បន្ទាត់ក្រិត ១:១ រ៉ាព័រទ័រ និងដែកឈាន \\ \hline
    \chkbox & កាត E4 (ទំព័រ ២), សន្លឹក ៣ តួនាទី E5 (ទំព័រ ៣) \& ផ្ទាំង \cnum{1}--\cnum{5} (ទំព័រ ៤) \\ \hline
    \chkbox & កាត Exit \textbf{A--D} (ទំព័រ ៥) \& ប័ណ្ណកិច្ចការផ្ទះ \#10, \#12 (ទំព័រ ៦) \\ \hline
  \end{tabular}
\end{minipage}%
\hfill
\begin{minipage}[t]{0.54\textwidth}
  \textbf{III.} {\muol កំហុសភ័ន្តច្រឡំ} \textbf{(តាមការស្រាវជ្រាវ) \&} {\muol ដំណោះស្រាយ}\\[2pt]
  \begin{tabular}{|L{9.25cm}|}
    \hline
    \rowcolor{headgray}
    \multicolumn{1}{|c|}{\textbf{ឧបសគ្គសិស្ស (Empirical Misconceptions) \& វិធានការទប់ស្កាត់}} \\ \hline
    \textbf{១. អន្ទាក់ ជ.ជ.ម (SSA) \& បន្ទាត់ជំនួយត្រឡប់រូប៖} ប្រៀបធៀបធ្នូ $35^\circ$ កាត់ ២ កន្លែង ($R_1, R_2$) ទល់កែង $90^\circ$ ($180^\circ$) និងគូស $AI$ ដាច់ៗ \en{(Yigit Koyunkaya, 2018; Muzaini et al., 2023)} ។ \\ \hline
    \textbf{២. ច្រឡំ Incenter vs Circumcenter/Centroid \& $IA=IB=IC$៖} បត់ក្រដាស $\tri{ABC}$ សាមញ្ញរក $r\approx 2.1\text{cm}$ ទល់ $IA, IB, IC$ \en{(Demiray et al., 2023; Boz-Yaman, 2024; Šipuš et al., 2022)} ។ \\ \hline
    \textbf{៣. រង្វង់តក្កទំ.១៥៤ \& សញ្ញា \textbf{/} លើទំ.១៥៥៖} កែទំ.១៥៤ (អ.ម), លុប \textbf{/} លើទំ.១៥៥, ប្រើ Flow-Chart ២ ដំណាក់កាល \en{(Anwar et al., 2021, 2023; Tchonang Youkap et al., 2020)} ។ \\ \hline
    \textbf{៤. ត្រីកោណត្រួតគ្នា (92\%$\to$53\%) \& គម្លាតទម្រង់ (85\% Level~0)៖} បំបែក $\tri{MDB}, \tri{MEC}$ / \#10 / \#12 + \textbf{[tracing]} + ៣ តួនាទី \en{(Biber, 2020; Manalo, 2025; Yeo \& Tzeng, 2019; Masinading \& Gaylo, 2022)} ។ \\ \hline
  \end{tabular}
\end{minipage}
}

\end{document}
